\section{Software design}

Dette afsnit beskriver den interne opbygning af systemet, samt de designprincipper og mønstre,
der er anvendt i implementeringen.
Hvor arkitekturafsnittet fokuserer på systemets overordnede struktur, fokuserer dette afsnit på den
konkrete realisering af komponenter og deres indbyrdes samarbejde.

\subsection{Designprincipper}

Systemet er designet med udgangspunkt i centrale software designprincipper med henblik på at sikre en
vedligeholdsbar og skalerbar kodebase.

\textbf{Separation of concerns:}
Ansvarsområder er opdelt mellem forskellige lag og komponenter.
Præsentationslaget håndterer udelukkende UI og brugerinteraktion, applikationslaget håndterer
forretningslogik, og datalaget står for persistens.

\textbf{Single Responsibility Principle (SRP):}
Hver komponent har ét veldefineret ansvar.
Eksempelvis håndterer \texttt{GraphQL} resolvers kun mapping mellem API-kald og service-laget,
mens services indeholder forretningslogik.

\textbf{Loose Coupling:}
Komponenter er designet med lav kobling, hvilket gør det muligt at ændre implementeringen af én
komponent uden at påvirke resten af systemet i væsentlig grad.

\textbf{DRY (Don't Repeat Yourself):}
Genbrug af kode er prioriteret gennem abstrahering af fælles funktionalitet i services og hjælpefunktioner.
Dette er også gjort brug af i præsentationslaget hvor der er implementeret forskellige components,
som kan genbruges, som \texttt{Button} eller \texttt{Card}.

\subsection{Intern struktur i applikationslaget}

Applikationslaget er opdelt i tre centrale lag: resolvers, services og dataadgang.

\textbf{Resolvers:}
\texttt{GraphQL} resolvers fungerer som indgangspunkt for klientforespørgsler.
De modtager queries og mutationer fra klienten og videresender disse til relevante services.
Resolvers indeholder minimal logik og fungerer primært som et forbindelsesled mellem API og
forretningslogik.

\textbf{Services:}
Service-laget indeholder systemets forretningslogik.
Her håndteres validering, databehandling og koordinering af kald til databasen.
Denne opdeling gør det muligt at genbruge logik på tværs af flere resolvers.

\textbf{Dataadgang:}
Dataadgang håndteres via Prisma, som fungerer som et abstraktionslag over PostgreSQL-databasen.
Prisma anvendes både til forespørgsler, migrationer og seeding af data.
Dette sikrer en konsistent og typesikker tilgang til databasen.

\subsection{Anvendte designmønstre}

Følgende designmønstre er anvendt i systemet:

\textbf{Layered Architecture Pattern:}
Systemet er struktureret i lag, hvor hvert lag har et klart defineret ansvar.
Dette forbedrer overskueligheden og vedligeholdbarheden.

\textbf{Resolver-Service Pattern:}
\texttt{GraphQL} resolvers er adskilt fra forretningslogikken, som er placeret i services.
Dette gør det muligt at teste logik uafhængigt af API-laget.

\textbf{Middleware Pattern:}
Middleware anvendes i \texttt{Express} til håndtering af tværgående funktionalitet som autentifikation,
logging og fejlbehandling.

\subsection{Dataflow og API-design}

Systemet anvender en kombination af \texttt{GraphQL} og \texttt{REST} til håndtering af klient-server
kommunikation.

Ved en typisk \texttt{GraphQL}-forespørgsel foregår dataflowet som følger:

\begin{enumerate}
    \item Klienten sender en query eller mutation via \texttt{Apollo Client}.
    \item Forespørgslen modtages af en \texttt{GraphQL} resolver.
    \item Resolveren kalder en relevant service.
    \item Servicen udfører nødvendig logik og interagerer med databasen via \texttt{Prisma}.
    \item Resultatet returneres gennem resolveren til klienten.
\end{enumerate}

Autentifikation håndteres separat via et \texttt{REST}-endpoint, hvor klienten sender login-oplysninger
gennem en \texttt{POST}-request.
Ved succesfuld autentifikation returneres en \texttt{JWT}, som efterfølgende anvendes i \texttt{GraphQL}-forespørgsler
til autorisation.

\subsection{Frontend design}

Frontend-applikationen er implementeret i \texttt{Next.js} og er opbygget omkring komponentbaseret udvikling.

UI'et består af genanvendelige komponenter, som hver især har et afgrænset ansvar.
State management håndteres primært gennem \texttt{Reacts useState} og Context API, hvilket er tilstrækkeligt,
givet systemets kompleksitet.

Datahentning foretages via \texttt{Apollo CLient}, som muliggør effektiv kommunikation med \texttt{GraphQL} API'et
samt caching af forespørgsler.

\subsection{Databasetilgang og persistens}

Databasen tilgås udelukkende gennem applikationslaget via \texttt{Prisma}.
Dette sikrer, at al adgang til data går gennem kontrollerede lag, hvor validering og forretningslogik
kan håndhæves.

\texttt{Prisma} anvendes til:
\begin{itemize}
    \item Databaseforespørgsler
    \item Migration af databasestruktur
    \item Seeding af initial data
\end{itemize}

Denne tilgang reducerer risikoen for inkonsistens og sikrer en struktureret udviklingsproces.

\subsection{Fejlhåndtering}

Fejlhåndtering er centraliseret gennem middleware i applikationslaget.
Dette sikrer ensartet håndtering af fejl og gør det muligt at returnere konsistente fejlmeddelelser
til klienten.

\texttt{GraphQL}-fejl håndteres i resolvers og services, hvor fejl enten håndteres lokalt eller videresendes
til central håndtering.

\subsection{Trade-offs}
I designet af systemet er der truffet en række bevidste valg, som indebærer forskellige trade-offs.
Valget af en monolitisk arkitektur reducerer kompleksiteten i både udvikling og deployment sammenlignet
med microservices.
Til gengæld kan det på sigt begrænse skalerbarheden, hvis systemet vokser betydeligt, hvor der
vil være brug for microservices.
Anvendelsen af \texttt{GraphQL} giver høj fleksibilitet i datahentning og reducerer over-fetching,
men introducerer samtidig øget kompleksitet i \texttt{API}'et, herunder håndtering af queries,
resolvers og performance optimering.
Kombinationen af \texttt{GraphQL} til datahåndtering og \texttt{REST} til autentifikation skaber en klar
ansvarsopdeling, men medfører også, at systemet anvender to forskellige \texttt{API}-paradigmer, hvilket
øger den samlede kompleksitet.
Endelig er der foretaget enkelte denormaliseringer i databasen for at forbedre performance, hvilket
sker på bekostning af streng normalisering og kan øge risikoen for inkonsistens, hvis det
ikke håndteres korrekt.

\subsection{Opsummering}

Softwaredesignet understøtter en modulær og vedligeholdbar struktur gennem klar ansvarsopdeling
og anvendelse af velkendte designprincipper og mønstre.
Kombination af \texttt{GraphQL}, \texttt{REST} og \texttt{Prisma} muliggør en fleksibel og effektiv
håndtering af data og forretningslogik.
Der er udarbejdet et C4 container diagram for at visualisere systemets overordnede struktur og interaktion
mellem centrale komponenter.